%% ----------------------------------------------------------------
%% Thesis.tex -- MAIN FILE (the one that you compile with LaTeX)
%% ---------------------------------------------------------------- 

% Set up the document
\documentclass[a4paper, 11pt, oneside]{Thesis}  % Use the "Thesis" style, based on the ECS Thesis style by Steve Gunn
\graphicspath{{Figures/}}  % Location of the graphics files (set up for graphics to be in PDF format)

% Include any extra LaTeX packages required
\usepackage[square, numbers, comma, sort&compress]{natbib}  % Use the "Natbib" style for the references in the Bibliography
\usepackage{verbatim}  % Needed for the "comment" environment to make LaTeX comments
\usepackage{vector}  % Allows "\bvec{}" and "\buvec{}" for "blackboard" style bold vectors in maths
\usepackage{textcomp} % Copyright symbol
\usepackage[]{algorithm2e}
\RestyleAlgo{algoruled}
\hypersetup{urlcolor=black, colorlinks=true}  % Colours hyperlinks in blue, but this can be distracting if there are many links.

%% ----------------------------------------------------------------
\begin{document}
\frontmatter	  % Begin Roman style (i, ii, iii, iv...) page numbering

% Set up the Title Page
\title  {Application of DiverseRank to Clustering}
%\authors  {\texorpdfstring
%{\href{your web site or email address}{Mayank Gupta}}
%{Mayank Gupta}
%}
\authors {Mayank Gupta\\201101004\\ \texttt{mayank.g@students.iiit.ac.in} \\ Ankush Jain\\201101010\\ \texttt{ankush.jain@students.iiit.ac.in}}
\addresses  {\groupname\\\deptname\\\univname}  % Do not change this here, instead these must be set in the "Thesis.cls" file, please look through it instead
\date       {\today}
\subject    {}
\keywords   {}

\maketitle
%% ----------------------------------------------------------------

\setstretch{1.3}  % It is better to have smaller font and larger line spacing than the other way round

% Define the page headers using the FancyHdr package and set up for one-sided printing
\fancyhead{}  % Clears all page headers and footers
\rhead{\thepage}  % Sets the right side header to show the page number
\lhead{}  % Clears the left side page header

\pagestyle{fancy}  % Finally, use the "fancy" page style to implement the FancyHdr headers

%% ----------------------------------------------------------------
% Copyright
\pagestyle{empty}  % No headers or footers for the following pages

\null\vfill
% Now comes the "Funny Quote", written in italics
%\textit{``Write a funny quote here.''}
\begin{center}
	Copyright {\textcopyright} Mayank Gupta, Ankush Jain, 2014 \\
	All Rights Reserved \\
\end{center}

%\begin{flushright}
%If the quote is taken from someone, their name goes here
%\end{flushright}

\vfill\vfill\vfill\vfill\vfill\vfill\null
\clearpage  % Funny Quote page ended, start a new page
%% ----------------------------------------------------------------
% Declaration Page required for the Thesis, your institution may give you a different text to place here
\Declaration{

	\addtocontents{toc}{\vspace{1em}}  % Add a gap in the Contents, for aesthetics

	\begin{center}
		It is certified thad the work contained in this project, titled \\
		\emph{"Application of DiverseRank to Classification"} \\ by \underline{Mayank Gupta} with Roll No. \underline{201101004} and \\
		\underline{Ankush Jain} with Roll No. \underline{201101010}\\
		has been carried out under my supervision for partial fulfillment of degree.\par
	\end{center}

	%I, AUTHOR NAME, declare that this thesis titled, `THESIS TITLE' and the work presented in it are my own. I confirm that:

	%\begin{itemize} 
	%\item[\tiny{$\blacksquare$}] This work was done wholly or mainly while in candidature for a research degree at this University.

	%\item[\tiny{$\blacksquare$}] Where any part of this thesis has previously been submitted for a degree or any other qualification at this University or any other institution, this has been clearly stated.

	%\item[\tiny{$\blacksquare$}] Where I have consulted the published work of others, this is always clearly attributed.

	%\item[\tiny{$\blacksquare$}] Where I have quoted from the work of others, the source is always given. With the exception of such quotations, this thesis is entirely my own work.

	%\item[\tiny{$\blacksquare$}] I have acknowledged all main sources of help.

	%\item[\tiny{$\blacksquare$}] Where the thesis is based on work done by myself jointly with others, I have made clear exactly what was done by others and what I have contributed myself.
	%\\
	%\end{itemize}


	Signed:\\
	\rule[1em]{25em}{0.5pt}  \\ % This prints a line for the signature
	(Adviser: Dr. P. Krishna Reddy)

	Date:\\
	\rule[1em]{25em}{0.5pt}  % This prints a line to write the date
}
\clearpage  % Declaration ended, now start a new page

%% ----------------------------------------------------------------
% The "Funny Quote Page"
\pagestyle{empty}  % No headers or footers for the following pages

\null\vfill
% Now comes the "Funny Quote", written in italics
\textit{``Winter is coming.''}

\begin{flushright}
	{\textemdash} Lord Eddard Stark
\end{flushright}

\vfill\vfill\vfill\vfill\vfill\vfill\null
\clearpage  % Funny Quote page ended, start a new page
%% ----------------------------------------------------------------

% The Abstract Page
\addtotoc{Abstract}  % Add the "Abstract" page entry to the Contents
\abstract{
	\addtocontents{toc}{\vspace{1em}}  % Add a gap in the Contents, for aesthetics

	The process of clustering divides a set of data objects into smaller sets consisting of similar objects.

	These clustering algorithms do not work efficiently in higher dimensional spaces because of the inherent sparsity of the data.

	Projected clustering algorithms have been proposed to find the clusters in hidden subspaces. The identification of more relevant attributes is a research issue in the projected clustering.

	In literature, an approach has been proposed to find the relevant subspaces using Frequent-Pattern Mining approaches.
	The frequent pattern can able to identify the item sets with high support.
	In frequent pattern based projected clustering approach, we observed that only frequent itemsets are considered.

	In this paper, we propose an approach to refine clusters generated by the frequent pattern based approach and produce more accurate clusters.
	In the proposed approach, we identify the irrelevant objects in a given cluster by exploiting the notion of diversity.
	Based on the notion of diversity, which captures the extent the items in the pattern belong to multiple categories of a concept hierarchy, the objects the membership of an object is identified.
	The membership of each object decides the existance of the object in the cluster.
	This process iteratively refines each cluster.
	We conduct the experiments on the real world datasets and show that the proposed approach improves the quality of the cluster

}

\clearpage  % Abstract ended, start a new page
%% ----------------------------------------------------------------

\setstretch{1.3}  % Reset the line-spacing to 1.3 for body text (if it has changed)

% The Acknowledgements page, for thanking everyone
\acknowledgements{
	\addtocontents{toc}{\vspace{1em}}  % Add a gap in the Contents, for aesthetics

	%The acknowledgements and the people to thank go here, don't forget to include your project advisor\ldots
	We are grateful to Mr. M. Kumaraswamy and Dr. P. Krishna Reddy for their constant help and guidance throughout the project.

}
\clearpage  % End of the Acknowledgements
%% ----------------------------------------------------------------

\pagestyle{fancy}  %The page style headers have been "empty" all this time, now use the "fancy" headers as defined before to bring them back


%% ----------------------------------------------------------------
\lhead{\emph{Contents}}  % Set the left side page header to "Contents"
\tableofcontents  % Write out the Table of Contents

%%% ----------------------------------------------------------------
%\lhead{\emph{List of Figures}}  % Set the left side page header to "List if Figures"
%\listoffigures  % Write out the List of Figures

%%% ----------------------------------------------------------------
%\lhead{\emph{List of Tables}}  % Set the left side page header to "List of Tables"
%\listoftables  % Write out the List of Tables

%%% ----------------------------------------------------------------
%\setstretch{1.5}  % Set the line spacing to 1.5, this makes the following tables easier to read
%\clearpage  % Start a new page
%\lhead{\emph{Abbreviations}}  % Set the left side page header to "Abbreviations"
%\listofsymbols{ll}  % Include a list of Abbreviations (a table of two columns)
%{
%% \textbf{Acronym} & \textbf{W}hat (it) \textbf{S}tands \textbf{F}or \\
%\textbf{LAH} & \textbf{L}ist \textbf{A}bbreviations \textbf{H}ere \\

%}

%%% ----------------------------------------------------------------
%\clearpage  % Start a new page
%\lhead{\emph{Physical Constants}}  % Set the left side page header to "Physical Constants"
%\listofconstants{lrcl}  % Include a list of Physical Constants (a four column table)
%{
%% Constant Name & Symbol & = & Constant Value (with units) \\
%Speed of Light & $c$ & $=$ & $2.997\ 924\ 58\times10^{8}\ \mbox{ms}^{-\mbox{s}}$ (exact)\\

%}

%%% ----------------------------------------------------------------
%\clearpage  %Start a new page
%\lhead{\emph{Symbols}}  % Set the left side page header to "Symbols"
%\listofnomenclature{lll}  % Include a list of Symbols (a three column table)
%{
%% symbol & name & unit \\
%$a$ & distance & m \\
%$P$ & power & W (Js$^{-1}$) \\
%& & \\ % Gap to separate the Roman symbols from the Greek
%$\omega$ & angular frequency & rads$^{-1}$ \\
%}
%%% ----------------------------------------------------------------
% End of the pre-able, contents and lists of things
% Begin the Dedication page

\setstretch{1.3}  % Return the line spacing back to 1.3

\pagestyle{empty}  % Page style needs to be empty for this page
\dedicatory{
	Dedicated to Shubham Sangal\ldots
}

\addtocontents{toc}{\vspace{2em}}  % Add a gap in the Contents, for aesthetics


%% ----------------------------------------------------------------
\mainmatter	  % Begin normal, numeric (1,2,3...) page numbering
\pagestyle{fancy}  % Return the page headers back to the "fancy" style

% Include the chapters of the thesis, as separate files
% Just uncomment the lines as you write the chapters

%% Chapter 1

\chapter{Chapter Title Here} % Write in your own chapter title
\label{Chapter0}
\lhead{Chapter 1. \emph{Chapter Title Here}} % Write in your own chapter title to set the page header


\section{Welcome and Thank You}

Welcome to this ``\LaTeX{} Thesis Template'', a beautiful and easy to use template for writing a thesis using the \LaTeX{} typesetting system.

If you are writing a thesis (or will be in the future) and its subject is technical or mathematical (though it doesn't have to be), then creating it in \LaTeX{} is highly recommended as a way to make sure you can just get down to the essential writing without having to worry over formatting or wasting time arguing with your word processor.

\LaTeX{} is easily able to professionally typeset documents that run to hundreds or thousands of pages long. With simple mark-up commands, it automatically sets out the table of contents, margins, page headers and footers and keeps the formatting consistent and beautiful. One of its main strengths is the way it can easily typeset mathematics, even \emph{heavy} mathematics. Even if those equations are the most horribly twisted and most difficult mathematical problems that can only be solved on a super-computer, you can at least count on \LaTeX{} to make them look stunning.

\subsection{Examples of \LaTeX{} Typeset eBooks}

You can see (and download and read) examples of eBooks typeset by \LaTeX{} in a small library here:\\
\href{http://www.sunilpatel.co.uk/typesetting.html}{\texttt{http://www.sunilpatel.co.uk/typesetting.html}}

The books are available for you to download for free. They were originally plain text files but were quickly and easily converted into beautifully typeset and well presented eBooks through \LaTeX{} for you to enjoy.


\section{Learning \LaTeX{}}

\LaTeX{} is not a WYSIWYG (What You See is What You Get) program, unlike word processors such as Microsoft Word or Corel WordPerfect. Instead, a document written for \LaTeX{} is actually a simple, plain text file that contains \emph{no formatting}. You tell \LaTeX{} how you want the formatting in the finished document by writing in simple commands amongst the text, for example, if I want to use \emph{italic text for emphasis}, I write the `$\backslash$\texttt{emph}\{\}' command and put the text I want in italics in between the curly braces. This means that \LaTeX{} is a ``mark-up'' language, very much like HTML.

\subsection{A (not so short) Introduction to \LaTeX{}}

If you are new to \LaTeX{}, there is a very good eBook -- freely available online as a PDF file -- called, ``The Not So Short Introduction to \LaTeX{}''. The book's title is typically shortened to just ``lshort''. You can download the latest version (as it is occasionally updated) from here:\\
\href{http://www.ctan.org/tex-archive/info/lshort/english/lshort.pdf}{\texttt{http://www.ctan.org/tex-archive/info/lshort/english/lshort.pdf}}

It is also available in several other languages. Find yours from the list on this page:\\
\href{http://www.ctan.org/tex-archive/info/lshort/}{\texttt{http://www.ctan.org/tex-archive/info/lshort/}}

It is recommended to take a little time out to learn how to use \LaTeX{} by creating several, small `test' documents. Making the effort now means you're not stuck learning the system when what you \emph{really} need to be doing is writing your thesis.

\subsection{A Short Math Guide for \LaTeX{}}

If you are writing a technical or mathematical thesis, then you may want to read the document by the AMS (American Mathematical Society) called, ``A Short Math Guide for \LaTeX{}''. It can be found online here:\\
\href{http://www.ams.org/tex/amslatex.html}{\texttt{http://www.ams.org/tex/amslatex.html}}\\
under the ``Additional Documentation'' section towards the bottom of the page.

\subsection{Common \LaTeX{} Math Symbols}
There are a multitude of mathematical symbols available for \LaTeX{} and it would take a great effort to learn the commands for them all. The most common ones you are likely to use are shown on this page:\\
\href{http://www.sunilpatel.co.uk/latexsymbols.html}{\texttt{http://www.sunilpatel.co.uk/latexsymbols.html}}

You can use this page as a reference or crib sheet, the symbols are rendered as large, high quality images so you can quickly find the \LaTeX{} command for the symbol you need.

\subsection{\LaTeX{} on a Mac}
 
The \LaTeX{} package is available for many systems including Windows, Linux and Mac OS X. The package for OS X is called MacTeX� and it contains all the applications you need � bundled together and pre-customised � for a fully working \LaTeX{} environment and work�ow.
 
MacTeX includes a dedicated \LaTeX{} IDE (Integrated Development Environment) called ``TeXShop'' for writing your `\texttt{.tex}' files and ``BibDesk'': a program to manage your references and create your bibliography section just as easily as managing songs and creating playlists in iTunes. To learn more about using \LaTeX{} on the Mac system, read the guide at:\\
\href{http://www.sunilpatel.co.uk/texhelp.html}{\texttt{http://www.sunilpatel.co.uk/texhelp.html }}


\section{Getting Started with this Template}

If you are familiar with \LaTeX{}, then you can familiarise yourself with the contents of the \href{http://www.sunilpatel.co.uk/files/Thesis\%20Template.zip}{Zip file} and the directory structure and then place your own information into the `\texttt{Thesis.cls}' file. Section \ref{FillingFile} on page \pageref{FillingFile} tells you how to do this. Make sure you read section \ref{ThesisConventions} about thesis conventions to get the most out of this template and then get started with the `\texttt{Thesis.tex}' file straightaway.

If you are new to \LaTeX{} it is recommended that you carry on reading through the rest of the information in this document.

\subsection{Latest Version}

This thesis template you have downloaded is version 1.2 and you can always find the latest version of this template online at:\\
\href{http://www.sunilpatel.co.uk/thesistemplate.html}{\texttt{http://www.sunilpatel.co.uk/thesistemplate.html}}

\subsection{List of Changes}

\subsubsection*{Version 1.2 (March 2008)}

This is version 1.2 and contains more corrections, improved advice on figures, an additional section for users on Mac systems and a major correction for those wishing to use US Letter sized paper for their thesis.

\subsubsection*{Version 1.1.2 (August 2007)}
This is version 1.1.2 and contains trivial (language and spelling) corrections.

\subsubsection*{Version 1.1.1 (May 2007)}
This is version 1.1.1 and corrects (updates) several hyperlinks that have changed. The file ``\texttt{lstpatch.sty}'' is also included as some \TeX{} distributions do not include it.

\subsubsection*{Version 1.1 (April 2007)}
This is version 1.1, released shortly after 1.0. It contains \emph{lots} of general text corrections (spelling, punctuation and grammar) and includes a link to a page containing a list of common \LaTeX{} math symbols, available at:\\
\href{http://www.sunilpatel.co.uk/latexsymbols.html}{\texttt{http://www.sunilpatel.co.uk/latexsymbols.html}}


\subsubsection*{Version 1.0 (March 2007)}
This is the first version, released to the public in March 2007.

\subsection{About this Template}

This \LaTeX{} Thesis Template is originally based and created around a \LaTeX{} style file created by Steve R.\ Gunn from the University of Southampton (UK), department of Electronics and Computer Science. You can find his original thesis style file at his site, here:\\
\href{http://www.ecs.soton.ac.uk/~srg/softwaretools/document/templates/}{\texttt{http://www.ecs.soton.ac.uk/$\sim$srg/softwaretools/document/templates/}}

My thesis originally used the `\texttt{ecsthesis.cls}' from his list of styles. However, I knew \LaTeX{} could still format better. To get the look I wanted, I modified his style and also created a skeleton framework and folder structure to place the thesis files in.

This Thesis Template consists of that modified style, the framework and the folder structure. All the work that has gone into the preparation and groundwork means that all you have to bother about is the writing.

Some of the modifications made to the original template include:
\begin{itemize}
\item Addition of the, `Declaration of Authorship' page
\item Optimised page layout for one-sided printing
\item Better page headers (changed from the default headers)
\item Full bibliography information (with `\texttt{natbib}' and `\texttt{unsrtnat}') that numerically sorts and compresses multiple references and also includes the URL field (with an active link) and DOI information if supplied
\item Change of font size and line spacing to increase readability
\item Added, `List of Symbols' page
\item Added, `funny quote' page
\end{itemize}


Before you begin using this template you should ensure that its style complies with the thesis style guidelines imposed by your institution. As this template is based on the style for the University of Southampton (UK), it can definitely be used without modification for Southampton PhD theses.

In most cases this template style and layout will be suitable. If it is not, it may only require a small change to bring the template in line with your institution's recommendations.


\section{What this Template Includes}

\subsection{Folders}

This template comes as a single \href{http://www.sunilpatel.co.uk/files/Thesis\%20Template.zip}{Zip file} that expands out to many files and folders. The folder names are mostly self-explanatory:

\textbf{Appendices} -- this is the folder where you put the appendices. Each appendix should go into its own separate `\texttt{.tex}' file.

\textbf{Chapters} -- this is the folder where you put the thesis chapters. A thesis usually has about seven chapters, though there is no hard rule on this. Each chapter should go in its own separate `\texttt{.tex}' file and they usually are split as:
\begin{itemize}
\item Chapter 1: Introduction to the thesis topic
\item Chapter 2: Background information and theory
\item Chapter 3: (Laboratory) experimental setup
\item Chapter 4: Details of experiment 1
\item Chapter 5: Details of experiment 2
\item Chapter 6: Discussion of the experimental results
\item Chapter 7: Conclusion and future directions
\end{itemize}
This chapter layout is specialised for the experimental sciences.

\textbf{Figures} -- this folder contains the all figures for the thesis. These are the final images that will go into the thesis document.

\textbf{Primitives} -- this is the folder that contains scraps, particularly because one final image in the `Figures' folder may be made from many separate images and photos, these source images go here. This keeps the intermediate files separate from the final thesis figures.

\subsection{Files}

Included are also several files, most of them are plain text and you can see their contents in a text editor. Luckily, many of them are auxiliary files created by \LaTeX{} or BibTeX and which you don't need to bother about:

\textbf{Bibliography.bib} -- this is an important file that contains all the bibliographic information and references that you will be citing in the thesis for use with BibTeX. You can write it manually, but there are reference manager programs available that will create and manage it for you. Bibliographies in \LaTeX{} are a large subject and you may need to read about BibTeX before starting with this.

\textbf{Thesis Todo List.txt} -- this is a Todo list that provides a simple way to track the thesis progress and also helps you see what parts need attention. Inside is a structured layout for you to enter all this information. Do not try to fill it out all at once, but write the tasks in as you come across them i.e.\ write an entry for a figure that needs creating only when you refer to or write about it (and hence need to place it) in the thesis text.

\textbf{Thesis.cls} -- this is an important file. It is the style file that tells \LaTeX{} how to format the thesis. You will also need to open this file in a text editor and fill in your own information (such as name, department, institution). Luckily, this is not too difficult and is explained in section \ref{FillingFile} on page \pageref{FillingFile}.

\textbf{Thesis.pdf} -- this is your beautifully typeset thesis (in the PDF file format) created by \LaTeX{}.

\textbf{Thesis.tex} -- this is an important file. This is the file that you tell \LaTeX{} to compile to produce your thesis as a PDF file. It contains the framework and constructs that tell \LaTeX{} how to layout the thesis. It is heavily commented so you can read exactly what each line of code does and why it is there. After you put your own information into the `\texttt{Thesis.cls}' file, go to this file and begin filling it in -- you have now started your thesis!

\textbf{vector.sty} -- this is a \LaTeX{} package, it tells \LaTeX{} how to typeset mathematical vectors. Using this package is very easy and you can read the documentation on the site (you just need to look at the `\texttt{vector.pdf}' file):\\
\href{http://www.ctan.org/tex-archive/macros/latex/contrib/vector/}{\texttt{http://www.ctan.org/tex-archive/macros/latex/contrib/vector/}}

\textbf{lstpatch.sty} -- this is a \LaTeX{} package required by this LaTeX template and is included as not all \TeX{} distributions have it installed by default. You do not need to modify this file.

Files that are \emph{not} included, but are created by \LaTeX{} as auxiliary files include:

\textbf{Thesis.aux} -- this is an auxiliary file generated by \LaTeX{}, if it is deleted \LaTeX{} simply regenerates it when you run the main `\texttt{.tex}' file.

\textbf{Thesis.bbl} -- this is an auxiliary file generated by BibTeX, if it is deleted, BibTeX simply regenerates it when you run the main tex file. Whereas the `\texttt{.bib}' file contains all the references you have, this `\texttt{.bbl}' file contains the references you have actually cited in the thesis and is used to build the bibliography section of the thesis.

\textbf{Thesis.blg} -- this is an auxiliary file generated by BibTeX, if it is deleted BibTeX simply regenerates it when you run the main `\texttt{.tex}' file.

\textbf{Thesis.lof} -- this is an auxiliary file generated by \LaTeX{}, if it is deleted \LaTeX{} simply regenerates it when you run the main `\texttt{.tex}' file. It tells \LaTeX{} how to build the `List of Figures' section.

\textbf{Thesis.log} -- this is an auxiliary file generated by \LaTeX{}, if it is deleted \LaTeX{} simply regenerates it when you run the main `\texttt{.tex}' file. It contains messages from \LaTeX{}, if you receive errors and warnings from \LaTeX{}, they will be in this `\texttt{.log}' file.

\textbf{Thesis.lot} -- this is an auxiliary file generated by \LaTeX{}, if it is deleted \LaTeX{} simply regenerates it when you run the main `\texttt{.tex}' file. It tells \LaTeX{} how to build the `List of Tables' section.

\textbf{Thesis.out} -- this is an auxiliary file generated by \LaTeX{}, if it is deleted \LaTeX{} simply regenerates it when you run the main `\texttt{.tex}' file.


So from this long list, only the files with the `\texttt{.sty}', `\texttt{.bib}', `\texttt{.cls}' and `\texttt{.tex}' extensions are the most important ones. The other auxiliary files can be ignored or deleted as \LaTeX{} and BibTeX will regenerate them.


\section{Filling in the `\texttt{Thesis.cls}' File}\label{FillingFile}

You will need to personalise the thesis template and make it your own by filling in your own information. This is done by editing the `\texttt{Thesis.cls}' file in a text editor.

Open the file and scroll down, past all the `$\backslash$\texttt{newcommand}\ldots' items until you see the entries for `\texttt{University Name}', `\texttt{Department Name}', etc\ldots.

Fill out the information about your group and institution and ensure you keep to block capitals where it asks you to. You can also insert web links, if you do, make sure you use the full URL, including the `\texttt{http://}' for this.

The last item you should need to fill in is the Faculty Name (in block capitals). When you have done this, save the file and recompile `\texttt{Thesis.tex}'. All the information you filled in should now be in the PDF, complete with web links. You can now begin your thesis proper!


\section{The `\texttt{Thesis.tex}' File Explained}

The \texttt{Thesis.tex} file contains the structure of the thesis. There are plenty of written comments that explain what pages, sections and formatting the \LaTeX{} code is creating. Initially there seems to be a lot of \LaTeX{} code, but this is all formatting, and it has all been taken care of so you don't have to do it.

Begin by filling out your information for the title page. For the thesis declaration, your institution may insist on something different than the text given. If this is the case, just replace what you see with what is required.

Then comes a page which contains a funny quote. You can put your own, or quote your favourite scientist, author, person, etc\ldots Make sure to put the name of the person who you took the quote from.

Next comes the acknowledgements. On this page, write about all the people who you wish to thank (not forgetting parents, partners and your advisor/supervisor).

The contents pages, list of figures and tables are all taken care of for you and do not need manually creating or editing. The next set of pages are optional and can be deleted since they are for a more technical thesis: Insert a list of abbreviations you have used in the thesis, then a list of the physical constants and numbers you refer to and finally, a list of mathematical symbols used in any formulae. Making the effort to fill these tables means the reader has a one-stop place to refer to instead of searching the internet and references to try and find out what you meant by certain abbreviations or symbols.

The list of symbols is split into the Roman and Greek alphabets. Whereas the abbreviations and symbols ought to be listed in alphabetical order (and this is \emph{not} done automatically for you) the list of physical constants should be grouped into similar themes.

The next page contains a one line dedication. Who will you dedicate your thesis to?

Finally, there is the section where the chapters are included. Uncomment the lines (delete the `\texttt{\%}' character) as you write the chapters. Each chapter should be written in its own file and put into the `Chapters' folder and named `\texttt{Chapter1}', `\texttt{Chapter2}, etc\ldots Similarly for the appendices, uncomment the lines as you need them. Each appendix should go into its own file and placed in the `Appendices' folder.

After the preamble, chapters and appendices finally comes the bibliography. The bibliography style (called `\texttt{unsrtnat}') is used for the bibliography and is a fully featured style that will even include links to where the referenced paper can be found online. Do not under estimate how grateful you reader will be to find that a reference to a paper is just a click away. Of course, this relies on you putting the URL information into the BibTeX file in the first place.

\section{Thesis Features and Conventions}\label{ThesisConventions}

To get the best out of this template, there are a few conventions that you may want to follow.

One of the most important (and most difficult) things to keep track of in such a long document as a thesis is consistency. Using certain conventions and ways of doing things (such as using the Todo list) makes the job easier. Of course, all of these are optional and you can adopt your own method.

\subsection{Printing Format}

This thesis template is designed for single sided printing as most theses are printed and bound this way. This means that the left margin is always wider than the right (for binding). Four out of �ve people will now judge the margins by eye and think, ``I never 
noticed that before.''.

The headers for the pages contain the page number on the right side (so it is easy to flick through to the page you want) and the chapter name on the left side.

The text is set to 11 point and a line spacing of 1.3. Generally, it is much more readable to have a smaller text size and wider gap between the lines than it is to have a larger text size and smaller gap. Again, you can tune the text size and spacing should you want or need to. The text size can be set in the options for the `$\backslash$\texttt{documentclass}' command at the top of the `\texttt{Thesis.tex}' file and the spacing can be changed by setting a different value in the `$\backslash$\texttt{setstretch}' commands (scattered throughout the `\texttt{Thesis.tex}' file).

\subsection{Using US Letter Paper}

The paper size used in the template is A4, which is a common � if not standard � size in Europe. If you are using this thesis template elsewhere and particularly in the United States, then you may have to change the A4 paper size to the US Letter size. Unfortunately, this is not as simple as replacing instances of `\texttt{a4paper}' with `\texttt{letterpaper}'.

This is because the final PDF file is created directly from the \LaTeX{} source using a program called `\texttt{pdfTeX}' and in certain conditions, paper size commands are ignored and all documents are created with the paper size set to the size stated in the configuration file for pdfTeX (called `\texttt{pdftex.cfg}').

What needs to be done is to change the paper size in the con�guration file for \texttt{pdfTeX} to reflect the letter size. There is an excellent tutorial on how to do this here: \\
\href{http://www.physics.wm.edu/~norman/latexhints/pdf_papersize.html}{\texttt{http://www.physics.wm.edu/$\sim$norman/latexhints/pdf\_papersize.html}}

It may be sufficient just to replace the dimensions of the A4 paper size with the US Letter size in the \texttt{pdftex.cfg} file. Due to the differences in the paper size, the resulting margins may be different to what you like or require (as it is common for Institutions to dictate certain margin sizes). If this is the case, then the margin sizes can be tweaked by opening up the \texttt{Thesis.cls} file and searching for the line beginning with, `$\backslash$\texttt{setmarginsrb}' (not very far down from the top), there you will see the margins speci�ed. Simply change those values to what you need (or what looks good) and save. Now your document should be set up for US Letter paper size with suitable margins.

\subsection{References}

The `\texttt{natbib}' package is used to format the bibliography and inserts references such as this one.\citep{Reference3} The options used in the `\texttt{Thesis.tex}' file mean that the references are listed in numerical order as they appear in the text. Multiple references are rearranged in numerical order\citep{Reference2, Reference1} and multiple, sequential references become reformatted to a reference range.\citep{Reference2, Reference1, Reference3} This is done automatically for you. To see how you use references, have a look at the `\texttt{Chapter1.tex}' source file.

References should come \emph{after} the punctuation mark if there is one (such as a comma or full stop). On the other hand, footnotes\footnote{Such as this footnote, here down at the bottom of the page.} come \emph{before} the punctuation mark. You can swap these around but the most important thing is to keep the convention consistent throughout the thesis. Footnotes themselves should be full, descriptive sentences (beginning with a capital letter and ending with a full stop).

To see how \LaTeX{} typesets the bibliography, have a look at the very end of this document (or just click on the reference number links).

\subsection{Figures}

There will hopefully be many figures in your thesis (that should be placed in the `Figures' folder). The way to insert figures into your thesis is to use a code template like this:
\begin{verbatim}
\begin{figure}[htbp]
  \centering
    \includegraphics{./Figures/Electron.pdf}
    \rule{35em}{0.5pt}
  \caption[An Electron]{An electron (artist's impression).}
  \label{fig:Electron}
\end{figure}
\end{verbatim}
Also look in the source file. Putting this code into the source file produces the picture of the electron that you can see in the figure below.

\begin{figure}[htbp]
	\centering
		\includegraphics{./Figures/Electron.pdf}
		\rule{35em}{0.5pt}
	\caption[An Electron]{An electron (artist's impression).}
	\label{fig:Electron}
\end{figure}

Sometimes figures don't always appear where you write them in the source. The placement depends on how much space there is on the page for the figure. Sometimes there is not enough room to fit a figure directly where it should go (in relation to the text) and so \LaTeX{} puts it at the top of the next page. Positioning figures is the job of \LaTeX{} and so you should only worry about making them look good!

Figures usually should have labels just in case you need to refer to them (such as in figure \ref{fig:Electron}). The `$\backslash$\texttt{caption}' command contains two parts, the first part, inside the square brackets is the title that will appear in the `List of Figures', and so should be short. The second part in the curly brackets should contain the longer and more descriptive caption text.

The `$\backslash$\texttt{rule}' command is optional and simply puts an aesthetic horizontal line below the image. If you do this for one image, do it for all of them.

The \LaTeX{} Thesis Template is able to use figures that are either in the PDF or JPEG file format. It is recommended that you read this short guide on how to get the best out of figures in \LaTeX{}, available here:\\
\href{http://www.sunilpatel.co.uk/texhelp5.html}{\texttt{http://www.sunilpatel.co.uk/texhelp5.html}}

Though it is geared more towards users of Mac and OS X systems, much of the advice applies to creating and using figures in general. It also explains why the PDF file format is preferred in figures over JPEG.

\subsection{Typesetting mathematics}

If your thesis is going to contain heavy mathematical content, be sure that \LaTeX{} will make it look beautiful, even though it won't be able to solve the equations for you.

The ``Not So Short Introduction to \LaTeX{}'' (available \href{http://www.ctan.org/tex-archive/info/lshort/english/lshort.pdf}{here}) should tell you everything you need to know for most cases of typesetting mathematics. If you need more information, a much more thorough mathematical guide is available from the AMS called, ``A Short Math Guide to \LaTeX{}'' and can be downloaded from:\\
\href{ftp://ftp.ams.org/pub/tex/doc/amsmath/short-math-guide.pdf}{\texttt{ftp://ftp.ams.org/pub/tex/doc/amsmath/short-math-guide.pdf}}

There are many different \LaTeX{} symbols to remember, luckily you can find the most common symbols \href{http://www.sunilpatel.co.uk/latexsymbols.html}{here}. You can use the web page as a quick reference or crib sheet and because the symbols are grouped and rendered as high quality images (each with a downloadable PDF), finding the symbol you need is quick and easy.

You can write an equation, which is automatically given an equation number by \LaTeX{} like this:
\begin{verbatim}
\begin{equation}
E = mc^{2}
  \label{eqn:Einstein}
\end{equation}
\end{verbatim}

This will produce Einstein's famous energy-matter equivalence equation:
\begin{equation}
E = mc^{2}
\label{eqn:Einstein}
\end{equation}

All equations you write (which are not in the middle of paragraph text) are automatically given equation numbers by \LaTeX{}. If you don't want a particular equation numbered, just put the command, `$\backslash$\texttt{nonumber}' immediately after the equation.


\section{Sectioning and Subsectioning}

You should break your thesis up into nice, bite-sized sections and subsections. \LaTeX{} automatically builds a table of Contents by looking at all the `$\backslash$\texttt{chapter}$\{\}$', `$\backslash$\texttt{section}$\{\}$' and `$\backslash$\texttt{subsection}$\{\}$' commands you write in the source.

The table of Contents should only list the sections to three (3) levels. A `$\backslash$\texttt{chapter}$\{\}$' is level one (1). A `$\backslash$\texttt{section}$\{\}$' is level two (2) and so a `$\backslash$\texttt{subsection}$\{\}$' is level three (3). In your thesis it is likely that you will even use a `$\backslash$\texttt{subsubsection}$\{\}$', which is level four (4). Adding all these will create an unnecessarily cluttered table of Contents and so you should use the `$\backslash$\texttt{subsubsection$^{*}\{\}$}' command instead (note the asterisk). The asterisk ($^{*}$) tells \LaTeX{} to omit listing the subsubsection in the Contents, keeping it clean and tidy.


\section{In Closing}

You have reached the end of this mini-guide. You can now rename or overwrite this pdf file and begin writing your own `\texttt{Chapter1.tex}' and the rest of your thesis. The easy work of setting up the structure and framework has been taken care of for you. It's now your job to fill it out! If you use this Thesis template and this mini-guide helps you, please let me know\footnote{Email: \href{mailto:web@sunilpatel.co.uk}{web@sunilpatel.co.uk} all comments and suggestions, questions and errata are welcome.}.

Good luck and have lots of fun!

\begin{flushright}
Sunil Patel: \href{http://www.sunilpatel.co.uk}{www.sunilpatel.co.uk}
\end{flushright}
 % Introduction

\chapter{Introduction}
\label{Chapter1}
\lhead{Chapter 1. \emph{Introduction}}

Clustering is a Data Mining technique, that, given a dataset, divides it into
groups of objects such that the objects in a group will be
similar to one another and different from the objects in other
groups\cite{bib17}. The objects are represented as a point, where
each dimension corresponds to an attribute/feature and the
feature value of each object determines its coefficient in the
corresponding dimension. The data points are grouped into
clusters based on the some similarity metric. The clustering
algorithms have large number of applications such as marketing (e.g., customer segmentation), image analysis, bio-
informatics, document classification, indexing, etc.

\section{Issues with Clustering Techniques}

Most clustering algorithms do not work efficiently in higher
dimensional space because of the inherent sparsity of the
data \cite{bib3}. In high dimensional space, it is likely that the
distance of any two points is almost the same for a large
class of common distributions \cite{bib5}. So, a clustering algorithm
employ feature selection based approach know as subspaces
clustering \cite{bib11}. The goal of the subspace clustering is to
find the particular dimensions on which the data points are
correlated and pruning away the remaining dimensions that
reduces the noise in the data.

The problem in these algorithms is that picking certain dimensions in advance can
lead to a loss of information. Furthermore, in many real
data, some points are correlated with respect to a given set
of dimensions and others are correlated with respect to different dimensions. Thus, it may not always be feasible to
prune off too many dimensions without considering the data
at the same time incurring a substantial loss of information.
Alternatively, the effects of dimensionality can be reduced
by a dimensionality reduction technique\cite{bib6}, however, in-
formation from all dimensions is uniformly transformed and
relevant information for some clusters may be reduced. Also,
the clusters may be hard to understand.

\section{Available Algorithms}

The projected clustering\cite{bib1} alogrithms overcome some of the
issues of the subspace clustering. In projected clustering, a
set of data points with an associated set of relevant dimensions are employed such that the data points are similar to
each other in the subspace formed by the relevant dimensions, but dissimilar to data point outside the cluster. The
widely used distance measures are more meaningful in projections of the high-dimensional space, where the object
values are dense\cite{bib9}. In other words, it is more likely for the data
to form dense, meaningful clusters in a high-dimensional
subspace. CLIQUE \cite{bib3}, PROCLUS \cite{bib1}, ORCLUS \cite{bib2}, and DOC \cite{bib12} are some efforts in the projected clustering approaches.

These approaches suffer from the quality of the clusters and
consumes lot of time. An effort is made to exploit the frequent pattern mining in the area of projected clustering to
address the issue of quality clusters\cite{bib18, bib19}. However, the
frequent pattern mining\cite{bib4} approach identifies the itemsets
(features/attributes) with high support. In frequent pattern
based projected clustering approach, we observed that only
frequent itemsets are considered and the infrequent itemsets
are completely eliminated from the clustering process.

\section{Our Approach}

We propose an approach to refine the clusters
generated by the frequent pattern based projected clustering approach and produce more qualitative clusters. We
observed that a high support pattern includes few items
and eliminates non-frequent items.
As a result of elimination of some of the items, the influence of other items
is not considered in the process of clustering. Due to elimination of some
data values, the the quality of clusters gets affected. In the
proposed approach, for a given cluster, we identify the mem-
bership of each data point employing the notion called diversity. The diversity identify the non-members of the given
cluster and remove the objects from the cluster. We conduct
the experiment on the real world datasets and show that the
proposed approach improves the quality of clusters.
 % Background Theory 

\chapter{Analysis Of Existing Algorithms}
\label{Chapter2}
\lhead{Chapter 2. \emph{Analysis Of Existing Algorithms}}

\section{CLIQUE, PROCLUS and ORCLUS}

One of the first algorithms in the area of projected clustering
algorithms is CLIQUE \cite{bib3}. CLIQUE finds the dense regions
(clusters) in a level-wise manner, based on the Apriori
principle. However, this algorithm does not scale well with
data dimensionality. In addition, the formed clusters have
large overlap, and this may not generate clear disjoint partitions. PROCLUS\cite{bib1} and ORCLUS\cite{bib2} employ alternative
techniques. They are much faster than CLIQUE and they
can discover disjoint clusters. In PROCLUS, the dimensions
relevant to each cluster are selected from the original set of
attributes. ORCLUS is more general and can select relevant attributes from the set of arbitrary directed orthogonal
vectors. PROCLUS fails to identify clusters with large difference in size and requires their dimensionality to be in a
predefined range. ORCLUS may discover clusters that are
hard to interpret.

\section{DOC}

DOC \cite{bib12} is a density-based algorithm that iteratively discovers projected clusters in a data set. DOC discovers one
cluster at a time. At each step, it tries to guess a good
medoid for the next cluster to be discovered. It repeatedly
picks a random point from the database and attempts to
discover the cluster centered at random point. Among all
discovered clusters, the cluster with the highest quality is selected. The process repeated for number of random points.

\section{Frequent Pattern-based Algorithm/MineClus}

The frequent pattern mining\cite{bib7} is one of the interesting
algorithms in data mining. The frequent itemset mining
problem was first proposed in\cite{bib4}. The frequent patterns
are generated from the transactional databases for a given
minsup. Subsequently, an efficient approach was proposed
in FP-growth\cite{bib8} to generate the frequent patterns without
generating the candidate items. In\cite{bib18, bib19}, a projected clustering algorithm was proposed. The algorithm iteratively
produce one cluster at a time by exploiting the frequent pattern mining approach. The top supported frequent patterns
are employed to find the item sets to generate the clusters.
A strength function is defined to arrange the clusters in to
strong and weak clusters. The strong clusters are produced
as an output and the objects in the weak clusters are added
by to the database. The process is iterated till all the objects
are grouped into the clusters. The FP-growth approach has
employed to generate the frequent patterns.
 % Experimental Setup

\chapter{MineClus and Diversity/DiverseRank}
\label{Chapter3}
\lhead{Chapter 3. \emph{MineClus and Diversity/DiverseRank}}

\section{MinClus}
The approach works based on the given a random medoid p $ \in $ S, the approach transform best projected cluster containing p, to the trasnactional databases. 
If the attribute of a record is bounded by p with respect to the width w (here, w=2), an item for that attribute is added to the corresponding itemset.
The appraoch observe that all frequent itemsets (i.e., combinations of dimensions) with respect to minsup = $\alpha * |S|$ are candidate clusters for medoid p.
The problem of finding the best projected cluster for a random medoid p can be transformed to the problem of finding the best itemset in a transformation of S, where goodness is defined by the $\mu$ function (refer Equation \ref{eq:1}).

Instead of discovering it in an non-deterministic way, a systematic data mining algorithm on S is applied. 
The association between the data items is identified by the frequent pattern mining algoirhtm. 
Due to association in the data items, the frequent pattern mining approach has been exploited in projected cluter approachs.
Recently, there is a more efficient algorithm, the FP-growth method \cite{bib8}. 
Here, the appraoch adopt frequent pattern mining for subspace clustering.
However, the objective is to find the frequent itemset with maximum $\mu$ value, rather than finding all frequent subspaces.
Assume that $(I_{best})$ is the itemset with the maximum $\mu$ value found so far and let $dim(I_{best})$ and $sup(I_{best})$ be its dimensionality and support, respectively. 

Let $I_{cond}$ be the current conditional pattern of the FP-growth process. Its support $sup(I_{cond})$ gives an upper bound for the supports of all patterns containing it. 
Moreover, the dimensionality of the itemsets that contain $I_{cond}$ is at most $dim(I_{cond})+l$, where l is the number of items above the items in in the header table of the FP-Growth. 
The $\mu$ function helps to find the strong and weak clusters. 
Among the strong clusters, the one with highest $\mu$ score is prouced as the final cluster and the data in the other clusters are added back to the database for generation of the next clusters. 
The process is repeated till the now data point is left in the database. 
The final data points which are not part of any cluster is treated as outliers.
\begin{equation}\label{eq:1}
\mu(a, b) = a * (1/\beta) * b
\end{equation}
where, a is support of the frequent patter, $\beta \in (0, 1]$ reflects
the importance of the projection, and b is the number of items in the frequent pattern.
%\figure{}

\section{Diversity/DiverseRank}
\subsection{About Concept Hierarchies}
A pattern contains data items. A concept hierarchy is a tree in which the data items are organized in a hierarchical manner. In this tree, all the leaf nodes are the items, the internal nodes represent the categories and the top node represents the root. The root could be a virtual node. Figure1 \iffalse(figure number) \fi is an example of the concept hierarchy.

Let C be a concept hierarchy. A node in C may be an item, category or root. The height of the root node is 0. Let n be a node in C. The height of n, is denoted as h(n), is equal to the number of edges on the path from root to n.

The concept hierarchies can be balanced and unbalanced. In the balanced concept hierarchy having height h has the same number of levels. The items at the given height are said to be at the same level. In an unbalanced concept hierarchy, the height of at least one of the leaf level node is different from the height of other leaf level nodes. The height of unbalanced concept hierarchy is equal to the height of the leaf level node having maximum height. In concept hierarchy C, all the lower-level nodes, except the root, are mapped to the immediate higher level nodes. We consider the concept hierarchies in which a lower level node is mapped to only one higher level node.

\subsection{Diversity}
The diversity of a pattern is based on the category of the items within it. If the items of a pattern are mapped to the same/few categories in a concept hierarchy, we consider that the pattern has low diversity. Relatively, if the items are mapped to multiple categories, we consider that the pattern has more diversity. 

\subsubsection{1}
For a given pattern diversity/diverserank is assigned based on the merging behavior in the corresponding concept hierarchy. If the pattern merges into few higher level categories quickly, it has low diversity/diverserank. Otherwise, if the pattern merges into one or a few high level categories slowly, it has relatively high diversity/diverserank value.

As an example, consider the concept hierarchy in Figure 1 \iffalse(figure number)\fi. For the pattern \{a, b\} the items a and b are mapped to the next level category p. In this case, the merging occurs quickly.For the pattern \{a, d\}, the item a is mapped to category p while item d is mapped to category q . Further, both th categories p and q are mapped to the category w. We say that the pattern \{a, d\} is more diverse than the pattern \{a, b\} as the merging is relatively slow in case of \{a, d\} as compared to \{a, b\}. Consider the pattern \{a, i\} which is relatively more diverse than the pattern \{a, d\} as both items merges at the root. The merging of \{a, i\} occurs slowly as comapred to \{a, d\}

\textbf{Computing the Diversity of Patterns}

We explain the process of calculating diverse rank of the pattern, called DRank, proposed in [\cite{bib14}, \cite{bib15}], the balanced and unbalanced pattern as follows.

We extract the projection of Concept Hierarchy for $Y$(pattern) to compute the diversity.

\textbf{Projection of Concept Hierarchy for $Y (P (Y /C))$}: Let $Y$ be $P$ and $C$ be concept hierarchy. The $P (Y /C)$ is the projection of $C$ for $Y$ which contains the portion of $C$. 
All the nodes and edges exists in the paths of the items of $Y$ to the root, along with the items and the root, are included in $P (Y /C)$. 
The projection $P (Y /C)$ is a tree which represents a concept hierarchy concerning to the pattern $Y$.

Given two patterns of the same length, different merging behavior can be realized, if we observe how the items in the patterns are mapped to higher level nodes. 
That is, one pattern may quickly merge to few higher level items within few levels and the other patterns may merge to few higher level items by crossing more number of levels. 
By capturing the process of merging, we define the notion of diverse rank (drank). 
So, $drank (Y )$ is calculated by capturing how the items are merged from leaf-level to root in $P (Y/C)$. 
It can be observed that a given item set maps from the leaf level to the root level through a merging process by crossing intermediate levels. 
At a given level, several lower level items/categories are merged into the corresponding higher level categories.

Two notions are employed to compute the diversity of pattern: Merging Factor ($MF$), Level Factor ($LF$) and Adjustment factor ($AF$).

We explain about MF after presenting the notion of generalized pattern.

\textbf{Generalized Pattern} $(GP (Y, l, P (Y /C)))$: Let $Y$ be a pattern, $h$ be the height of $P (Y /C)$ and $l$ be an integer.
The $GP (Y, l, P (Y /C))$ indicates the $GP$ of $Y$ at level  $l$ in $P (Y /C)$. 
Assume that the $GP (Y, l + 1, P (Y /C))$ is given. 
The $GP (Y, l, P (Y /C))$ is calculated based on the $GP$ of $Y$ at level $(l + 1)$. 
The $GP (Y, l, P (Y /C))$ is obtained by replacing every item at level $(l + 1)$ in $GP (Y, l + 1, P (Y /C))$ with its corresponding parent at the level $l$ with duplicates removed, if any.

The notion of merging factor at level l is defined as follows.

\textbf{Merging factor} $(MF (Y, l, P (Y /C)))$: Let $Y$ be pattern and $l$ be the level. The merging factor indicates how the items of a
pattern merge from the level $l + 1$ to the level $l (0 ≤ l < h)$.
If there is no change, the $MF(Y,l)$ is $1$. 
If all items merges to one node, the $MF(X,l)$ value equals to $0$.
So, the $MF$ value at the level $l$ is denoted by $MF(Y, l, P (Y /C))$ which is equal
to the ratio of the number of nodes in $(GP (Y, l, P (Y /C) − 1)$ to the number of nodes in $(GP (Y, l + 1, P (Y /C) − 1)$.

\begin{equation}
    MF(Y,l,P(Y/C)) = \frac{|GP(Y,l,P(Y/C)| - 1}{|GP(Y,l+1,P(Y/C)))|-1}
\end{equation}

We now define the notion of level factor to determine the contribution of nodes at the given level.

Level Factor $(LF (l, P (Y /C))$: For a given $P (Y /C)$, $h$ be the height of $P (Y /C) = {0, 1}$. 
Let $l$ be such that $1 ≤ l ≤ (h − 1)$.
The $LF$ value of $P (Y /C)$ at level $l$ indicates the contribution of nodes at level $l$ to $DRank$.
We can assign equal, linear or exponential weights to each level. 
Here, we provide a formula which assigns the weight to the level such that the weight is in proportion the level number.

\begin{equation}
LF (l, P (Y /C))=\frac{2 * (h - l)}{ h * (h - 1)}
\end{equation}

\textbf{Adjustment factor} $AF (Y, l, P (Y /C))$:Let $Y$ be pattern and $l$ be the level
The Adjustment Factor $(AF)$ at level $l$ helps in reducing the drank by measuring the contribution of dummy edges/nodes relative to the original edges/nodes at the level $l$. 
The $AF$ for a pattern $Y$ at a level $l$ should depend on the ratio of number of real edges formed with the children of the 
real nodes in $P (Y /E)$ versus total number of edges formed with the children of real and dummy nodes at $l$ in $P (Y /E)$. 
The value of $AF$ at a given height should lie between $0$ and $1$. 
If the number of real edges is equals to zero, $AF$ is $0$.
If the pattern at the given level does not contain dummy nodes/edges, the value of $AF$ becomes $1$.
Note that the $AF$ value is not defined at the leaf level nodes as children do not exist. 
The $AF$ for $Y$ at level $l$ is denoted as $AF (Y, l, P (Y /E))$ and is calculated by the following formula.
The value for $AF$ is less than $1$ when the pattern is unbalanced otherwise it returns $1$.

\begin{equation}
    AF (Y, l, P (Y /C)) = \frac{\# of Real Edges of U P (Y, l, P (Y /E)) }{ \# of T otal Edges of U P (Y, l, P (Y /E))}
\end{equation}

where numerator is the number of edges formed with the children of the real nodes and denominator is the number of edges formed with the children of both real and dummy nodes at the level $l$ in $P (Y /E)$.

The approach to compute drank of $Y$ is as follows. 
We convert the concept hierarchy to extended concept hierarchy called, "extended concept hierarchy" by adding dummy nodes and edges. 
Next, we adjust the drank in accordance with the number of dummy nodes and edges using the notion called adjustment factor. 
So, the drank of $Y$ is relative to the drank of the same pattern computed by considering all of its items are at the leaf level of the extended concept hierarchy. 
The dummy nodes and edges are added when the pattern is unbalanced.

\textbf{Diverse rank of a frequent pattern} $Y (drank(Y))$: Let $Y$ be the pattern and $C$ be the unbalanced concept hierarchy
of height $h$. 
The drank of $Y$ , denoted by $drank(Y )$, is given by the following equation:

\begin{equation}
    drank(Y, C) = [M F (Y, l, P (Y /E)) ∗ AF (Y, l, P (Y /E))] * LF (l, P (Y /E))     
\end{equation}

where, $h$ is the height of the $P (P/E)$, $E$ is the extended unbalanced concept hierarchy, $MF (Y, l, P (Y /E))$ is the $MF$ 
of $Y$ at level $l$, $LF (l, P (Y /E))$ is the $LF$ at level $l$ and $AF (Y, l, P (Y /E))$ is the $AF$ of $Y$ at level $l$.

\subsubsection{2}
For a given pattern diversity/diverserank is assigned based on difference in number of edges in the sub tree of corresponding concept hierarchy which only contain those leaves(items) which are present in the pattern (e') and minimum number of edges needed to form a tree of height same as the corresponding concept hierarchy  number of leaf nodes is equal to number of distinct items in the pattern(e). 
To normalize it we divide it with the difference of number of edges if all the pattern merges at root(E) and  minimum number of edges needed to form a tree of height same as the     corresponding concept hierarchy  number of leaf nodes is equal to number of distinct items in the pattern(e)

\begin{equation}\label{eq:2}
    dr(Y,C)=\frac{e'-e}{E-e}
\end{equation}

As an example, consider the concept hierarchy in Figure 1 \iffalse(figure number)\fi. \\ 
For the pattern \{a,b\}, E=6, e'=4 and e=4. So diversty=0.0 \\
For the pattern \{a,d\}, E=6, e'=5 and e=4. So diversty=0.5 \\
For the pattern \{a,i\}, E=6, e'=6 and e=4. So diversty=1.0
 % Experiment 1

%\chapter{Proposed Approach}
\label{Chapter4}
\lhead{Chapter 4. \emph{Proposed Approach}}

Our approach uses Diversity to refine each cluster as it is formed. The MineClus algorithm generates clusters iteratively using a \emph{Frequent Pattern} with high support.
Right after a cluster is formed, we try to validate the membership of each element in the newly formed cluster. Any element that is deemed to be different from the rest of the cluster is put back into the dataset, and the MineClus algorithm will consider it while forming newer clusters in the future. Hence, an element is added to the cluster which is best suited for it. If a suitable cluster cannot be found for the element, it is listed as a singleton cluster.

The rationale and key idea behind using diversity is that all elements of the cluster will have similar diversity values with respect to the Concept Hierarchy. The concept hierarchy captures domain knowledge of the dataset and represents it in a tree-like form.
Diversity is able to identify how different a data point
is from other points in the cluster. If the data point
is too different (exceeding a threshold), it clearly does not belong to the current cluster and is labelled as a non-member.
All non-members are put back into the list of data points that are yet to be clustered.

For the purposes of refining, let us consider a cluster \(C\) containing \(n\) data points \(\{o_1, o_2, \dots, o_n \}\). We use diversity to determine the membership of \(o_i\). The array containing the diversity values for each of the cluster elements is denoted as \(\{d_1, d_2, \dots, d_n\}\). Two approaches for this are proposed.

\section{Diverse Square-Root Distance}

We define a measure called \emph{Diverse Square Root Distance (dsrd(k))} as:

\begin{equation}
dsrd(k) = \sqrt{\sum_{j=1}^{n} {(d_k - d_j)}^2}
\end{equation}

\(dsrd(k)\) represents the distance of the \(k^{th}\) element from the rest of the cluster.

\begin{algorithm}[H]
 \KwData{\(C, D, \delta\)}
$C$ is the current cluster\;
$D$ contains elements to be clustered\;
$\delta$ is a user-defined threshold\;
 \KwResult{Refined $C$}
 %\While{not at end of this document}{
\For{$i \leftarrow 1$ to $n$}{
  %read current\;
  %\eIf{understand}{
   %go to next section\;
   %current section becomes this one\;
   %}{
   %go back to the beginning of current section\;
  %}
\eIf{$dsrd(i) > \delta$} {
	Remove $i^{th}$ element from $C$\;
}{
	Keep $i^{th}$ element in $C$\;
}
 }
return $C$\;
 \caption{How to write algorithms}
\end{algorithm}
 % Experiment 2

\chapter{Experiments}
\label{Chapter5}
\lhead{Chapter 5. \emph{Experiments}}

\section{Work Done}
\begin{itemize}
	\item We studied relevant literature, as mentioned in the references
	\item We understood and implemented the Apriori Algorithm and FP-Growth algorithms in Python
	\item We understood and implemented both the implementions of Diversity in Python
	\item We implemented the entire MinClus clustering suite with support for merging
	\item We proposed and implemented six Outlier detection/Cluster refining methods
	\item We conducted extensive testing of our approach and that of the original MinClus approach, results of which are presented below.
\end{itemize}

\section{Evaluation Metric}
We view Clustering as a series of decisions - for each pair of items in the dataset, whether to put them in the same cluster or in a different cluster.

A true positive (TP) decision assigns two
similar documents to the same cluster, a true negative (TN)
decision assigns two dissimilar documents to different clus-
ters. There are two types of errors we can commit. A (FP)
decision assigns two dissimilar documents to the same clus-
ter. A (FN) decision assigns two similar documents to dif-
ferent clusters.

\section{Results with Diversity1(original approach)}
\subsection{Without merging}

\begin{table}[H]
\caption{Accuracy} 
\begin{center}
		\begin{tabular}{ | l | l | l | l |}
				\hline

				\textbf{Dataset} & \textbf{MinClus} & \textbf{DSRD} & \textbf{Binning} \\ \hline
                \textbf{Iris} & 82.20 & 85.58 & 85.58 \\ \hline
                \textbf{Seed} & 79.11 & 78.96 & 78.96 \\ \hline
                \textbf{Zoo} & 88.33 & 88.55 & 81.18  \\ \hline
                \textbf{Water} & 60.23 & 62.28 & 60.24  \\ \hline
		\end{tabular}
\end{center}
\label{table:acc1}
\end{table}



\begin{table}[H]
\caption{Precision} 
\begin{center}
		\begin{tabular}{ | l | l | l | l |}
				\hline

				\textbf{Dataset} & \textbf{MinClus} & \textbf{DSRD} & \textbf{Binning} \\ \hline

				\textbf{Iris} & 97.50 & 98.28 & 98.28 \\ \hline
				\textbf{Seed} & 90.62 & 89.03 & 89.03 \\ \hline
				\textbf{Zoo} & 94.02 & 91.28 & 88.94  \\ \hline
				\textbf{Water} & 36.32 & 39.07 & 36.60 \\ \hline
		\end{tabular}
\end{center}
\label{table:pre1}
\end{table}


\begin{table}[H]
\caption{Recall} 
\begin{center}
		\begin{tabular}{ | l | l | l | l |}
				\hline

				\textbf{Dataset} & \textbf{MinClus} & \textbf{DSRD} & \textbf{Binning} \\ \hline

				\textbf{Iris} &  47.84 & 57.76 & 57.76 \\ \hline
				\textbf{Seed} & 41.65 & 42.05 & 42.05 \\ \hline
				\textbf{Zoo} & 55.03 & 57.96 & 24.89  \\ \hline
				\textbf{Water} & 19.61 & 15.79 & 20.09 \\ \hline
		\end{tabular}
\end{center}
\label{table:rec1}
\end{table}

\subsection{With merging}

\begin{table}[H]
\caption{Accuracy} 
\begin{center}
		\begin{tabular}{ | l | l | l | l |}
				\hline

				\textbf{Dataset} & \textbf{MinClus} & \textbf{DSRD} & \textbf{Binning} \\ \hline
                \textbf{Iris} & 82.78 & 91.80 & 91.80 \\ \hline
                \textbf{Seed} & 82.51 & 83.21 & 83.21 \\ \hline
                \textbf{Zoo} & 91.02 & 94.65 & 90.47  \\ \hline
                \textbf{Water} & 60.29 & 59.80 & 60.50  \\ \hline
		\end{tabular}
\end{center}
\label{table:acc2}
\end{table}



\begin{table}[H]
\caption{Precision} 
\begin{center}
		\begin{tabular}{ | l | l | l | l |}
				\hline

				\textbf{Dataset} & \textbf{MinClus} & \textbf{DSRD} & \textbf{Binning} \\ \hline

				\textbf{Iris} & 75.89 & 87.38 & 87.38 \\ \hline
				\textbf{Seed} & 78.72 & 79.86 & 79.86 \\ \hline
				\textbf{Zoo}  & 90.95 & 90.02 & 81.06 \\ \hline
				\textbf{Water}& 36.94 & 36.70 & 39.38 \\ \hline
		\end{tabular}
\end{center}
\label{table:pre2}
\end{table}

\begin{table}[H]
\caption{Recall} 
\begin{center}
		\begin{tabular}{ | l | l | l | l |}
				\hline

				\textbf{Dataset} & \textbf{MinClus} & \textbf{DSRD} & \textbf{Binning} \\ \hline

				\textbf{Iris}  & 70.85 & 88.11 & 88.11 \\ \hline
				\textbf{Seed}  & 65.13 & 66.35 & 66.35  \\ \hline
				\textbf{Zoo}   & 69.61 & 87.45 & 78.82  \\ \hline
				\textbf{Water} & 20.63 & 22.06 & 25.90 \\ \hline
		\end{tabular}
\end{center}
\label{table:rec2}
\end{table}


\section{Results with Diversity2(improved approach)}
\subsection{Without merging}

\begin{table}[H]
\caption{Accuracy} 
\begin{center}
		\begin{tabular}{ | l | l | l | l |}
				\hline

				\textbf{Dataset} & \textbf{MinClus} & \textbf{DSRD} & \textbf{Binning} \\ \hline
                \textbf{Iris} & 82.20 & 82.43 & 77.54 \\ \hline
                \textbf{Seed} & 79.11 & 79.96 & 75.06 \\ \hline
                \textbf{Zoo} & 88.33 & 88.33 & 85.65  \\ \hline
                \textbf{Water} & 60.23 & 63.35 & 60.27  \\ \hline
		\end{tabular}
\end{center}
\label{table:acc1}
\end{table}



\begin{table}[H]
\caption{Precision} 
\begin{center}
		\begin{tabular}{ | l | l | l | l |}
				\hline

				\textbf{Dataset} & \textbf{MinClus} & \textbf{DSRD} & \textbf{Binning} \\ \hline

                \textbf{Iris} & 97.50 & 98.74 & 98.34 \\ \hline
                \textbf{Seed} & 90.62 & 95.76 & 91.00 \\ \hline
                \textbf{Zoo} & 94.02 & 94.02 & 96.79  \\ \hline
                \textbf{Water} & 36.32 & 39.96 & 36.37 \\ \hline
		\end{tabular}
\end{center}
\label{table:pre1}
\end{table}


\begin{table}[H]
\caption{Recall} 
\begin{center}
		\begin{tabular}{ | l | l | l | l |}
				\hline

				\textbf{Dataset} & \textbf{MinClus} & \textbf{DSRD} & \textbf{Binning} \\ \hline

                \textbf{Iris} &  47.84 & 47.89 & 33.17 \\ \hline
                \textbf{Seed} & 41.65 & 41.74 & 27.93 \\ \hline
                \textbf{Zoo} & 55.03 & 55.03 & 41.75  \\ \hline
                \textbf{Water} & 19.61 & 11.43 & 19.53 \\ \hline
		\end{tabular}
\end{center}
\label{table:rec1}
\end{table}

\subsection{With merging}

\begin{table}[H]
\caption{Accuracy} 
\begin{center}
		\begin{tabular}{ | l | l | l | l |}
				\hline

				\textbf{Dataset} & \textbf{MinClus} & \textbf{DSRD} & \textbf{Binning} \\ \hline
                \textbf{Iris} & 82.78 & 90.94 & 93.24 \\ \hline
                \textbf{Seed} & 82.51 & 85.00 & 83.41 \\ \hline
                \textbf{Zoo} & 91.02 & 91.96 & 95.78  \\ \hline
                \textbf{Water} & 60.29 & 62.05 & 60.18  \\ \hline
		\end{tabular}
\end{center}
\label{table:acc2}
\end{table}



\begin{table}[H]
\caption{Precision} 
\begin{center}
		\begin{tabular}{ | l | l | l | l |}
				\hline

				\textbf{Dataset} & \textbf{MinClus} & \textbf{DSRD} & \textbf{Binning} \\ \hline

                \textbf{Iris} & 75.89 & 86.21 & 89.82 \\ \hline
                \textbf{Seed} & 78.72 & 77.50 & 80.45 \\ \hline
                \textbf{Zoo}  & 90.95 & 86.58 & 93.10 \\ \hline
                \textbf{Water}& 36.94 & 38.86 & 36.71 \\ \hline
		\end{tabular}
\end{center}
\label{table:pre2}
\end{table}

\begin{table}[H]
\caption{Recall} 
\begin{center}
		\begin{tabular}{ | l | l | l | l |}
				\hline

				\textbf{Dataset} & \textbf{MinClus} & \textbf{DSRD} & \textbf{Binning} \\ \hline

                \textbf{Iris}  & 70.85 & 86.69 & 89.92 \\ \hline
                \textbf{Seed}  & 65.13 & 77.51 & 66.38 \\ \hline
                \textbf{Zoo}   & 69.61 & 78.82 & 89.08 \\ \hline
                \textbf{Water} & 20.63 & 16.57 & 20.59 \\ \hline
		\end{tabular}
\end{center}
\label{table:rec2}
\end{table}

 % Results and Discussion

%\input{./Chapters/Chapter7} % Conclusion

%% ----------------------------------------------------------------
% Now begin the Appendices, including them as separate files

\addtocontents{toc}{\vspace{2em}} % Add a gap in the Contents, for aesthetics

\appendix % Cue to tell LaTeX that the following 'chapters' are Appendices

% Appendix A

\chapter{Appendix Title Here}
\label{AppendixA}
\lhead{Appendix A. \emph{Appendix Title Here}}

Write your Appendix content here.	% Appendix Title

%\input{./Appendices/AppendixB} % Appendix Title

%\input{./Appendices/AppendixC} % Appendix Title

\addtocontents{toc}{\vspace{2em}}  % Add a gap in the Contents, for aesthetics
\backmatter

%% ----------------------------------------------------------------
\label{Bibliography}
\lhead{\emph{Bibliography}}  % Change the left side page header to "Bibliography"
\bibliographystyle{unsrtnat}  % Use the "unsrtnat" BibTeX style for formatting the Bibliography
%\bibliography{Bibliography}  % The references (bibliography) information are stored in the file named "Bibliography.bib"
\begin{thebibliography}{99}
	\bibitem{bib1}
		C. C. Aggarwal, J. L. Wolf, P. S. Yu, C. Procopiuc, and J. S. Park. Fast algorithms for projected clustering. \emph{SIGMOD} Rec., 28(2):61–72, June 1999.
	\bibitem{bib2}
		C. C. Aggarwal and P. S. Yu. Finding generalized projected clusters in high dimensional spaces. \emph{SIGMOD} Rec., 29(2):70–81, May 2000.
	\bibitem{bib3}
		R. Agrawal, J. Gehrke, D. Gunopulos, and P. Raghavan. Automatic subspace clustering of high dimensional data for data mining applications.  SIGMOD Rec., 27(2):94–105, June 1998.  
	\bibitem{bib4}
		R. Agrawal and R. Srikant. Fast algorithms for mining association rules in large databases. \emph{In Proceedings of the 20th International Conference on Very Large Data Bases}, VLDB ’94, pages 487–499, San Francisco, CA, USA, 1994. Morgan Kaufmann Publishers Inc.
	\bibitem{bib5}
		K. S. Beyer, J. Goldstein, R. Ramakrishnan, and U. Shaft. When is “nearest neighbor” meaningful? \emph{In Proceedings of the 7th International Conference on Database Theory}, ICDT ’99, pages 217–235, London, UK, UK, 1999. Springer-Verlag.
	\bibitem{bib6}
		C. Faloutsos and K.-I. Lin. Fastmap: A fast algorithm for indexing, data-mining and visualization of traditional and multimedia datasets. \emph{SIGMOD} Rec., 24(2):163–174, May 1995.
	\bibitem{bib7}
		J. Han, H. Cheng, D. Xin, and X. Yan. Frequent
		pattern mining: Current status and future directions.
		\emph{Data Min. Knowl. Discov.}, 15(1):55–86, 2007.
	\bibitem{bib8}
		J. Han, J. Pei, and Y. Yin. Mining frequent patterns
		without candidate generation. \emph{SIGMOD} Rec.,
		29(2):1–12, May 2000.
	\bibitem{bib9}
		A. Hinneburg, C. C. Aggarwal, and D. A. Keim. What
		is the nearest neighbor in high dimensional spaces? \emph{In
			Proceedings of the 26th International Conference on
		Very Large Data Bases}, VLDB ’00, pages 506–515,
		San Francisco, CA, USA, 2000. Morgan Kaufmann
		Publishers Inc.
	\bibitem{bib10}
		A. K. Jain, M. N. Murty, and P. J. Flynn. Data
		clustering: A review. \emph{ACM Comput. Surv.},
		31(3):264–323, 1999.
	\bibitem{bib11}
		G. Moise, A. Zimek, P. Kruger, H.-P. Kriegel,
		and J. Sander. Subspace and projected clustering:
		Experimental evaluation and analysis. \emph{Knowl. Inf.
		Syst.}, 21(3):299–326, 2009.
	\bibitem{bib12}
		C. M. Procopiuc, M. Jones, P. K. Agarwal, and T. M.
		Murali. A monte carlo algorithm for fast projective
		clustering. In \emph{Proceedings of the 2002 ACM SIGMOD
		International Conference on Management of Data},
		SIGMOD ’02, pages 418–427, New York, NY, USA,
		2002. ACM.
	\bibitem{bib13}
		W. M. Rand. Objective criteria for the evaluation of
		clustering methods. \emph{Journal of the American
		Statistical Association}, 66(336):846–850, 1971.
	\bibitem{bib14}
		S. Srivastava, R. U. Kiran, and P. K. Reddy.
		Discovering diverse-frequent patterns in transactional
		databases. In \emph{Proceedings of the 17th International
		Conference on Management of Data}, COMAD ’11,
		pages 14:1–14:10. Computer Society of India, 2011.
	\bibitem{bib15}
		M. K. Swamy, P. K. Reddy, and S. Srivastava.
		Extracting diverse patterns with unbalanced concept
		hierarchy. In \emph{Advances in Knowledge Discovery and
		Data Mining - 18th Pacific-Asia Conference}, PAKDD
		2014, Tainan, Taiwan, May 13-16, 2014. Proceedings,
		Part I, pages 15–27. Springer-Verlag, 2014.
	\bibitem{bib16}
		UCI. Machine Learning Repository,
		\url{https://archive.ics.uci.edu/ml/index.html}, 2014
		(accessed August, 2014).
	\bibitem{bib17}
		R. Xu and D. Wunsch, II. Survey of clustering
		algorithms. \emph{Trans. Neur. Netw.}, 16(3):645–678, 2005.
	\bibitem{bib18}
		M. L. Yiu and N. Mamoulis. Frequent-pattern based
		iterative projected clustering. In \emph{Proceedings of the
			Third IEEE International Conference on Data
		Mining}, ICDM ’03, pages 689–692, Washington, DC,
		USA, 2003. IEEE Computer Society.
	\bibitem{bib19}
		M. L. Yiu and N. Mamoulis. Iterative projected
		clustering by subspace mining. \emph{IEEE Trans. on
		Knowl. and Data Eng.}, 17(2):176–189, 2005.
\end{thebibliography}

\end{document}  % The End
%% ----------------------------------------------------------------
